% GENERATED FILE. Edit canonical lessons, then run npm run generate:books.
% canonical-source-hash: fnv1a64:4dac042fba0548be
% canonical-lessons: RU-C413-proyezd, RU-C413-platskart, RU-C413-kupe, RU-C413-tarif, RU-C413-lgota

\chapter{A train ticket}
\label{ch:ru-a-train-ticket}
\hlchaptermodality{413}

\begin{chapteropening}
\textbf{By the end of this chapter:} \emph{I can buy a train ticket: the fare, an open sleeping car, a compartment, the rate and a concession.}

Complete the last lesson of chapter 413: I can buy a train ticket: the fare, an open sleeping car, a compartment, the rate and a concession.
\end{chapteropening}

\section[\texorpdfstring{\ru{проезд} (proyézd)}{проезд (proyézd)}]{\ru{проезд} (proyézd) — a fare; a passage, a way through}

\label{lesson:RU-C413-proyezd}



\begin{quote}
Before the new one: say the Russian for to let through, then the Russian for road markings.
\end{quote}

\subsection*{You'll want to know: \ru{проезд}}
\textbf{\ru{проезд}} (\emph{proyézd}) — \textquotedblleft{}a fare; a passage, a way through\textquotedblright{}.

\begin{quote}
\ru{Сколько} \ru{стоит} \ru{проезд} \ru{на} \ru{метро}?
\end{quote}

\begin{quote}
\emph{Skólko stóit proyézd na metró?}
\end{quote}

\begin{quote}
How much is the metro fare?
\end{quote}

Masculine. \emph{\ru{Проезд} \ru{закрыт}} — no through road.

\begin{cousinweb}
\emph{\ru{проезд}} is built from \emph{\ru{про}-} 'through' and \emph{-\ru{езд}} 'riding', the root of \emph{\ru{ездить}} 'to ride, to travel'. EtymDB lists \emph{\ru{проезд}} without a deeper source.
\end{cousinweb}

\begin{grammarlens}[title={What you've built}]
One more word for telling, explaining and describing at B1.
\end{grammarlens}

\subsection*{Your turn}
\begin{itemize}
\raggedright
  \item \emph{Say it:} \emph{proyézd}
  \item \emph{Say it:} \emph{proyézd}, once more
  \item \emph{Say it:} a sentence of your own with \emph{proyézd}
  \item \emph{Recall:} say the Russian for to let through, then the Russian for road markings, then say \emph{proyézd} again
\end{itemize}

\begin{tcolorbox}[breakable,colback=teal!4,colframe=teal!35!black,title={Before you move on}]
What does \ru{проезд} mean? (\textquotedblleft{}A fare; a passage, a way through\textquotedblright{}.) Say it once more.
\end{tcolorbox}
\section[\texorpdfstring{\ru{плацкарт} (platskárt)}{плацкарт (platskárt)}]{\ru{плацкарт} (platskárt) — an open sleeping car (on a train)}

\label{lesson:RU-C413-platskart}



\begin{quote}
Before the new one: say the Russian for road markings, then the Russian for a fare.
\end{quote}

\subsection*{You'll want to know: \ru{плацкарт}}
\textbf{\ru{плацкарт}} (\emph{platskárt}) — \textquotedblleft{}an open sleeping car (on a train)\textquotedblright{}.

\begin{quote}
\ru{Студенты} \ru{ехали} \ru{на} \ru{море} \ru{в} \ru{плацкарте}.
\end{quote}

\begin{quote}
\emph{Studénty yékhali na móre v platskárte.}
\end{quote}

\begin{quote}
The students travelled to the seaside in an open sleeping car.
\end{quote}

Masculine; colloquial for \emph{\ru{плацкартный} \ru{вагон}}, a cheap carriage of open bunks.

\begin{cousinweb}
\emph{\ru{плацкарт}} is usually traced to German \emph{Platzkarte} 'a seat reservation', built from \emph{Platz} 'place' and \emph{Karte} 'card'. EtymDB lists the German word but not the Russian, so the route is hedged.
\end{cousinweb}

\begin{grammarlens}[title={What you've built}]
One more word for telling, explaining and describing at B1.
\end{grammarlens}

\subsection*{Your turn}
\begin{itemize}
\raggedright
  \item \emph{Say it:} \emph{platskárt}
  \item \emph{Say it:} \emph{platskárt}, once more
  \item \emph{Say it:} a sentence of your own with \emph{platskárt}
  \item \emph{Recall:} say the Russian for road markings, then the Russian for a fare, then say \emph{platskárt} again
\end{itemize}

\begin{tcolorbox}[breakable,colback=teal!4,colframe=teal!35!black,title={Before you move on}]
What does \ru{плацкарт} mean? (\textquotedblleft{}An open sleeping car (on a train)\textquotedblright{}.) Say it once more.
\end{tcolorbox}
\section[\texorpdfstring{\ru{купе} (kupé)}{купе (kupé)}]{\ru{купе} (kupé) — a compartment (on a train)}

\label{lesson:RU-C413-kupe}



\begin{quote}
Before the new one: say the Russian for a fare, then the Russian for an open sleeping car.
\end{quote}

\subsection*{You'll want to know: \ru{купе}}
\textbf{\ru{купе}} (\emph{kupé}) — \textquotedblleft{}a compartment (on a train)\textquotedblright{}.

\begin{quote}
\ru{В} \ru{купе} \ru{четыре} \ru{полки}.
\end{quote}

\begin{quote}
\emph{V kupé chetýre pólki.}
\end{quote}

\begin{quote}
There are four berths in a compartment.
\end{quote}

Neuter and indeclinable: \emph{\ru{в} \ru{купе}}, \emph{\ru{из} \ru{купе}}.

\begin{cousinweb}
\emph{\ru{купе}} is usually traced to French \emph{coupé} 'cut off', from \emph{couper} 'to cut': a carriage cut into small rooms. EtymDB does not list the Russian word, so the line is hedged.
\end{cousinweb}

\begin{grammarlens}[title={What you've built}]
One more word for telling, explaining and describing at B1.
\end{grammarlens}

\subsection*{Your turn}
\begin{itemize}
\raggedright
  \item \emph{Say it:} \emph{kupé}
  \item \emph{Say it:} \emph{kupé}, once more
  \item \emph{Say it:} a sentence of your own with \emph{kupé}
  \item \emph{Recall:} say the Russian for a fare, then the Russian for an open sleeping car, then say \emph{kupé} again
\end{itemize}

\begin{tcolorbox}[breakable,colback=teal!4,colframe=teal!35!black,title={Before you move on}]
What does \ru{купе} mean? (\textquotedblleft{}A compartment (on a train)\textquotedblright{}.) Say it once more.
\end{tcolorbox}
\section[\texorpdfstring{\ru{тариф} (taríf)}{тариф (taríf)}]{\ru{тариф} (taríf) — a rate, a tariff; a fare}

\label{lesson:RU-C413-tarif}



\begin{quote}
Before the new one: say the Russian for an open sleeping car, then the Russian for a compartment.
\end{quote}

\subsection*{You'll want to know: \ru{тариф}}
\textbf{\ru{тариф}} (\emph{taríf}) — \textquotedblleft{}a rate, a tariff; a fare\textquotedblright{}.

\begin{quote}
\ru{С} \ru{января} \ru{в} \ru{автобусе} \ru{новый} \ru{тариф}.
\end{quote}

\begin{quote}
\emph{S yanvaryá v avtóbuse nóvyy taríf.}
\end{quote}

\begin{quote}
From January there is a new fare on the bus.
\end{quote}

Masculine. \emph{\ru{Тариф}} is the official rate; \emph{\ru{цена}} the everyday word for a price.

\begin{cousinweb}
\emph{\ru{тариф}} was borrowed from Italian \emph{tariffa} 'a tariff', which EtymDB traces to Arabic \emph{taʿrīfa} 'a rate'.
\end{cousinweb}

\begin{grammarlens}[title={What you've built}]
One more word for telling, explaining and describing at B1.
\end{grammarlens}

\subsection*{Your turn}
\begin{itemize}
\raggedright
  \item \emph{Say it:} \emph{taríf}
  \item \emph{Say it:} \emph{taríf}, once more
  \item \emph{Say it:} a sentence of your own with \emph{taríf}
  \item \emph{Recall:} say the Russian for an open sleeping car, then the Russian for a compartment, then say \emph{taríf} again
\end{itemize}

\begin{tcolorbox}[breakable,colback=teal!4,colframe=teal!35!black,title={Before you move on}]
What does \ru{тариф} mean? (\textquotedblleft{}A rate, a tariff; a fare\textquotedblright{}.) Say it once more.
\end{tcolorbox}
\section[\texorpdfstring{\ru{льгота} (lgóta)}{льгота (lgóta)}]{\ru{льгота} (lgóta) — a concession, a reduced fare; a benefit}

\label{lesson:RU-C413-lgota}



\begin{quote}
Before the new one: say the Russian for a compartment, then the Russian for a rate.
\end{quote}

\subsection*{You'll want to know: \ru{льгота}}
\textbf{\ru{льгота}} (\emph{lgóta}) — \textquotedblleft{}a concession, a reduced fare; a benefit\textquotedblright{}.

\begin{quote}
\ru{У} \ru{пенсионеров} \ru{есть} \ru{льгота} \ru{на} \ru{проезд}.
\end{quote}

\begin{quote}
\emph{U pensionérov yest lgóta na proyézd.}
\end{quote}

\begin{quote}
Pensioners have a reduced fare on public transport.
\end{quote}

Feminine. \emph{\ru{Льготный} \ru{билет}} — a concession ticket.

\begin{cousinweb}
\emph{\ru{льгота}} is usually linked with an old root meaning 'light, easy', the one in \emph{\ru{лёгкий}}: a concession makes things lighter. EtymDB lists the word without a source, so the line is hedged.
\end{cousinweb}

\begin{grammarlens}[title={What you've built}]
One more word for telling, explaining and describing at B1.
\end{grammarlens}

\subsection*{Your turn}
\begin{itemize}
\raggedright
  \item \emph{Say it:} \emph{lgóta}
  \item \emph{Say it:} \emph{lgóta}, once more
  \item \emph{Say it:} a sentence of your own with \emph{lgóta}
  \item \emph{Recall:} say the Russian for a compartment, then the Russian for a rate, then say \emph{lgóta} again
\end{itemize}

\begin{tcolorbox}[breakable,colback=teal!4,colframe=teal!35!black,title={Before you move on}]
What does \ru{льгота} mean? (\textquotedblleft{}A concession, a reduced fare; a benefit\textquotedblright{}.) Say it once more.
\end{tcolorbox}
